
\documentclass[../../main.tex]{subfiles}
\begin{document}

% 年度の大きな表示
\begin{center}
  {\fontsize{48pt}{58pt}\selectfont \textbf{1979}\normalsize\textbf{年度}}
\end{center}

\vspace{10mm}

% 本文

\vspace{15mm}

% 出題テーマと難易度の表
\noindent\textbf{出題テーマと難易度}

\vspace{5mm}

\begin{center}
  \shadowbox{%
    \begin{tabular}{|c|p{8cm}|c|}
      \hline
      \textbf{問題番号} & \textbf{テーマ} & \textbf{難易度} \\
      \hline
      第1問 & 空間ベクトルと極限 & \\
      \hline
      第2問 & 三角関数の最大値 & \\
      \hline
      第3問 & 積分方程式 & \\
      \hline
      第4問 & 対数関数と面積比 & \\
      \hline
    \end{tabular}%
  }
\end{center}

\vspace{15mm}

% 解答
\noindent\textbf{解答}

\vspace{5mm}

\begin{center}
  \shadowbox{%
    \renewcommand{\arraystretch}{1.6}
    \begin{tabular}{|c|c|p{8cm}|}
      \hline
      \textbf{問題番号} & \textbf{小問} & \textbf{答え} \\
      \hline
      第1問 & - & $P_n\to\left(\dfrac{1}{13},\dfrac{1}{13},\dfrac{1}{13}\right)$，$Q_n\to\left(-\dfrac{5}{13},\dfrac{11}{26},\dfrac{5}{26}\right)$ \\
      \hline
      第2問 & - & $-1\le a\le8$ \\
      \hline
      \multirow{2}{*}{第3問} & (1) & 証明問題 \\
      \cline{2-3}
                             & (2) & $P(x)=\dfrac{1}{3}x^3+x^2+3x,\ Q(x)=\dfrac{1}{3}x^3+x^2+5x+2$ \\
      \hline
      第4問 & - & $S:T=\{c(\log c-1)-a(\log a-1)\}:\dfrac{1}{2}(c\log c-a\log a)$ \\
      \hline
    \end{tabular}%
    \renewcommand{\arraystretch}{1.0}
  }
\end{center}

\vspace{15mm}

% 解答時間
\noindent\textbf{試験形態}

\vspace{5mm}

\begin{center}
  \shadowbox{%
    \begin{tabular}{p{8cm}}
      （要確認）
    \end{tabular}%
  }
\end{center}

\end{document}
